\section*{Was Version 1.6 verändert und korrigiert}
\addcontentsline{toc}{section}{Neu in v1.6: Ergebnisse, Reichweite und Lesepfad}
Diese Fassung vom 15. September 2026 führt das vollständige Hauptdokument v1.5 fort. Die 25 bisherigen Kapiteldateien einschließlich aller darin enthaltenen Hauptkapitel und Anhänge sowie sämtliche Abbildungen bleiben erhalten. Dreizehn inhaltliche Einfügungen verbinden die neuen Ergebnisse direkt mit den bisherigen Argumenten. Fünf neue Hauptkapitel und ein Evidenzanhang führen die Rechnungen, Modellgrenzen und Arbeitsaufträge zusammen. Historische Formulierungen werden als historische Stände gelesen, nicht als aktuell offene oder aktuell erledigte Aufgaben.

\begin{bild}{Der wichtigste neue Zusammenhang}
Der Baukasten ist kleiner geworden: Eine konkrete Ordnung aus komplexen $2\times2$-Matrizen trägt tatsächlich das markierte $E_8$-Gitter. Aus denselben internen Matrizen lassen sich die Amplituden eines Weyl-Walks bilden. Damit sind zwei vorher getrennt erzählte Baupläne ausdrücklich verbunden. Es fehlen aber weiterhin die aus der Quelle erzeugten Ortsverschiebungen und die Auswahl von Energie, Anfangszustand und Aufzeichnungen. Ein gemeinsames Alphabet ist jetzt genauer bekannt; die physisch ausgewählte Ausführung ist noch nicht bestimmt.
\end{bild}

Der Matrixbeweis ist \emph{keine neue Entdeckung dieser Revision}: Er steht bereits in einer Repository-Untersuchung vom 9. September. Der letzte Status hatte dessen vollständigen Abschnitt und die markierte Compiler-Abbildung nicht berücksichtigt. v1.6 korrigiert dieses Versäumnis und reproduziert den vorhandenen Prüfer. Neu hinzu kommen die direkte Rechnung der Walk-Koeffizienten im selben Matrixkörper und ein unabhängig nachgerechneter Voll-Fock-Zustandsvergleich.

Die drei jüngsten Eingaben werden nicht pauschal bestätigt. Bestätigt beziehungsweise präzisiert sind Matrixordnung, $A_3$-Gewichtsgeometrie, signierte native Kopplung und Teile der Fock-Analyse. Nicht nachgewiesen bleiben die behauptete Ableitung der Raumdimension allein aus $\mu_4$, eine automatische Dreifamilienauswahl, die Entfernung der Spiegelmoden durch Faltung und die Gleichsetzung eines inneren Wurzellabels mit einer physischen Position.

Außerdem werden erstmals in dieser PDF-Linie die bereits als vorläufige Markdown-Texte vorliegenden v1.6-Ergebnisse vollständig eingeordnet: Ein-Record-Präparation, unbedingte Vier-Record-Wahrscheinlichkeiten, Spektralzertifikat des benannten C16-Singulettmodells, endliche Transportrechnungen und der gemeinsame kosmologische Parametertest. Diese älteren großen Läufe wurden in dieser Revision nicht sämtlich erneut ausgeführt. Ihre übernommene Evidenz wird von den fünf frisch wiederholten Prüfprogrammen getrennt.

\textbf{Lesepfad.} Kapitel~\ref{ch:v16-labor} konsolidiert die noch nicht ins PDF übernommenen Resultate; Kapitel~\ref{ch:v16-matrix} erklärt die Matrixbrücke; Kapitel~\ref{ch:v16-operations} trennt innere Operation und Bewegung. Kapitel~\ref{ch:v16-state} behandelt die Zustandsauswahl, Kapitel~\ref{ch:v16-program} die sechs Fragen und T1--T8. Anhang~\ref{app:v16-evidence} nennt Quellen, Prüfungen und Grenzen. Das separate Update v1.6 ist keine gekürzte Ersatzfassung des Buches.

\begin{grenze}{Keine globale Schließung aus lokalen Nachweisen}
Die neuen Prüfungen schließen weder T1--T8 noch RH, Faktorisierung oder P versus NP. Sie liefern eine präzisere gemeinsame algebraische Darstellung, nachprüfbare positive Teilergebnisse und konkrete Gegenmodelle gegen zu starke Herkunftsbehauptungen. Ein dynamischer Spin-2-Sektor und ein gemeinsames chirales 3+1D-Maß sind nicht konstruiert.
\end{grenze}
